% chapters/02_system.tex

\chapter{System im Überblick}
\label{chap:system}
\label{chap:grundlagen}

\Cref{fig:architektur} zeigt den Weg vom Datensatz über Training und Checkpoint bis zum
Policy-Server, den die Simulation abfragt.

\begin{figure}[htb]
  \centering
  \resizebox{\linewidth}{!}{%
  \begin{tikzpicture}[
    node distance = 0.7cm and 2.0cm,
    box/.style   = {draw, rounded corners, fill=NavyBlue!8, align=center,
                    minimum width=3.0cm, minimum height=1.25cm, font=\small\sffamily},
    ext/.style   = {box, fill=OliveGreen!10},
    miss/.style  = {box, fill=white, dashed, draw=gray},
    arrow/.style = {-{Stealth[scale=1.1]}, thick, NavyBlue!70},
    lab/.style   = {font=\scriptsize\sffamily, align=center},
  ]
    \node[ext] (ds)   {Datensatz (\acs{hf})\\\scriptsize LeRobot v3.0, 301 Ep.};
    \node[box] (tr)   [right=of ds] {Training\\\scriptsize Docker lokal /\\\scriptsize Apptainer KISSKI};
    \node[ext] (ck)   [right=of tr] {Checkpoint (\acs{hf})\\\scriptsize \texttt{luca-mue/}\\\scriptsize \texttt{groot-g1dex3-*}};
    \node[box] (srv)  [below=1.3cm of ck] {Policy-Server\\\scriptsize GR00T, ZMQ :5555};
    \node[box] (sim)  [left=of srv] {Sim-Eval\\\scriptsize Isaac Lab, 30\,Hz};
    \node[miss] (rob) [left=of sim] {Roboter G1\\\scriptsize nicht angebunden};

    \draw[arrow] (ds) -- node[above, lab] {Download,\\v3.0$\to$v2.1} (tr);
    \draw[arrow] (tr) -- node[above, lab] {Upload} (ck);
    \draw[arrow] (ck) -- node[right, lab] {laden} (srv);
    \draw[arrow, <->] (sim) -- node[below, lab] {Bilder, Zustand /\\16$\times$28 Ziele} (srv);
    \draw[arrow, <->, dashed, gray] (rob.south) -- ++(0,-0.5) -| node[pos=0.25, below, lab, gray] {gleiches Protokoll} (srv.south);
  \end{tikzpicture}}
  \caption{Datenfluss des Systems. Grün: externe Ablagen auf Hugging Face; gestrichelt: nicht
    umgesetzt. Formate und Ports: \cref{chap:schnittstellen}.}
  \label{fig:architektur}
\end{figure}

\section{Plattform}
\label{sec:plattform}
\label{sec:humanoide-robotik}

Der Unitree~G1~\cite{unitree_g1} wird mit der dreifingrigen DEX3-Hand verwendet. Gesteuert
werden nur Arme und Hände: je sieben Gelenke pro Arm und pro Hand, zusammen 28~\acp{dof}.
Beine und Taille bleiben in der Simulation fest. Die vier Kameras des Datensatzes (zwei am
Kopf, eine an jedem Handgelenk) liefern Bilder mit 640$\times$480 Pixeln bei 30~Hz.

\section{GR00T N1.6 und N1.7}
\label{sec:vla}
\label{sec:imitation-learning}
\label{sec:flow-matching}

GR00T gehört wie OpenVLA~\cite{openvla2024} und $\pi_0$~\cite{pi0_2024} zu den \acp{vla} und verbindet ein Vision-Language-Modell mit einem Aktionskopf, der per Flow
Matching~\cite{lipman2023flow} einen Block von 16 Gelenkbefehlen erzeugt~\cite{groot2025}.
Nachtrainiert wird per \acf{bc}~\cite{hussein2017imitation}: Das Modell lernt, zu jedem
beobachteten Zustand die demonstrierte Aktion vorherzusagen. Der Trainings-Loss misst daher
nur die Nachahmung, nicht den Aufgabenerfolg; weicht die Policy im Regelkreis von den
Demonstrationen ab, summieren sich Fehler~\cite{ross2011dagger}.

\textbf{N1.6}~\cite{groot_n16} (3,29~Mrd.\ Parameter, Backbone Eagle, Python~3.10) ist der Standardpfad; verwendet wird ein Fork des NVIDIA-Repositorys~\cite{groot_repo}.
\textbf{N1.7}~\cite{groot_n17} (3,14~Mrd.\ Parameter) nutzt das Backbone Cosmos-Reason2-2B (Qwen3-VL, zugangsbeschränkt) und läuft in
einem eigenen Image (Python~3.12). Checkpoints beider Generationen sind nicht gegenseitig
ladbar; die Wahl erfolgt über \texttt{GROOT\_VERSION}. N1.7 wurde mit den NVIDIA-Standardwerten
trainiert, die vom N1.6-Rezept abweichen (u.\,a.\ State-Dropout 0,2 statt 0); der Vergleich
beider ist deshalb einer von Modell \emph{und} Rezept
(\doc{docs/weiterfuehrend/groot-n17-migration.md}).

\section{Rechenumgebungen}
\label{sec:rechenumgebungen}
\label{sec:infrastruktur}
\label{sec:kisski}

\begin{table}[htb]
  \centering
  \begin{tabularx}{\linewidth}{lXl}
    \toprule
    Umgebung & Zweck & GPU \\
    \midrule
    KISSKI (GWDG)~\cite{kisski} & Training, Open-Loop-Auswertung (\ac{slurm}, Apptainer, max.\ 48~h) & 4$\times$A100 80~GB, H100 94~GB \\
    \acs{ikr}-Server (Institut) & Simulation, visuelles Gate, Rendering & 2$\times$RTX PRO 6000 \\
    lokal / vast.ai & kleine Trainings, Sim-Eval auf Mietkarten & je nach Mietkarte \\
    \bottomrule
  \end{tabularx}
  \caption[Rechenumgebungen]{Rechenumgebungen. Isaac~Sim braucht Raytracing-Kerne (\acs{rt}-Cores); A100 und H100 haben keine, die
    Simulation läuft daher nie auf KISSKI.}
  \label{tab:rechenumgebungen}
\end{table}

Training und Simulation laufen in Containern (\acs{cuda}~12.8, \texttt{uv}), konfiguriert nur über
Umgebungsvariablen. Checkpoints wandern über Hugging Face vom Cluster zum Sim-Server.
